\documentclass[10pt,a4paper]{amsart}
\usepackage[margin=19mm]{geometry}
\usepackage{amsmath,amssymb,mathtools}
\usepackage[scr=rsfs,cal=euler]{mathalfa}
\usepackage{xfrac}
\usepackage[unicode,hidelinks,bookmarks=false]{hyperref}
\newcommand{\ie}{i.e.\ }
\newcommand{\cf}{cf.\ }
\newcommand{\resp}{respectively\ }
\newcommand{\Subsection}{\subsection}
\providecommand{\nobreakdash}{\nobreak\hbox{-}}
\setlength{\emergencystretch}{2em}
\multlinegap0pt
\usepackage{bm}
\usepackage{bbm}
\usepackage{dsfont}
\usepackage{xspace}
\usepackage{cancel}
\usepackage{mathtools}
\usepackage{amsthm}
\makeatletter
\@ifundefined{theorem}{\newtheorem{theorem}{Theorem}}{}
\@ifundefined{definition}{\newtheorem{definition}{Definition}}{}
\@ifundefined{lemma}{\newtheorem{lemma}[theorem]{Lemma}}{}
\makeatother
\title[A diagrammatic approach to the three-page index]{A diagrammatic approach to the three-page index}
\author{Hyungkee Yoo}
\date{Source: arXiv:2601.12846v1 (CC BY 4.0)}
\begin{document}
\maketitle

\noindent\emph{Source and licence.} A diagrammatic approach to the three-page index. Version arXiv:2601.12846v1 (CC BY 4.0), distributed under the Creative Commons Attribution 4.0 International licence, as recorded in the arXiv licence metadata for this exact version and shown on the abstract page https://arxiv.org/abs/2601.12846v1. Full rights notice: licenses/arxiv-2601.12846-CC-BY-4.0.txt.\par\smallskip

\noindent\textit{Editorial scope.} Counted unit: the Lemma whose source label is \texttt{lem:extended} in the article (source0.tex lines 275--279, source label \texttt{lem:extended}), delivered together with the article's own complete original proof of it (source0.tex lines 281--321, one \texttt{proof} environment of 41 lines). No mathematics is written, repaired or re-derived: statement and proof are the article's own text line for line. Required context retained uncounted: def:binding-circle (lines 167--181); lem:binding-circle-existence (lines 189--192). Every definition, display and label of those lines is retained unchanged. Omitted from this package and not redistributed: 1--166, 182--188, 193--274, 280--280, 322--507. Nothing inside a retained excerpt is omitted or altered: every retained body is delivered exactly as the article's lines are written, so no omission receipt of any kind accompanies this package. Citations: the retained excerpts carry no citation command at all, so no bibliography entry of the article is transcribed and no reference is redistributed. Reference adaptation: none was needed and none was made; every reference inside the delivered unit resolves inside it, so no unresolved reference or citation is produced. Printed slips kept literal: no printed word, symbol, hypothesis or number of the counted statement or of its proof was altered. Layout and class: the article's own class is \texttt{amsart} with packages including graphicx, multirow, array, amsmath, amssymb, xcolor, longtable, supertabular; the wrapper is the lane's pinned \texttt{amsart} editorial wrapper at 10pt on A4, which restates the theorem-like environments the retained text needs and adds the article's own macro definitions that the retained text needs (none), with extra wrapper packages (bm, bbm, dsfont, xspace, cancel, mathtools). The delivered numbering is the unit's own. Assets: no retained span contains a figure environment, a graphics-inclusion command, a PDF-page inclusion, a table or a TikZ picture; no image file is copied into this package, the package declares no asset dependency and no image file is redistributed with it. Licence: version arXiv:2601.12846v1 of this article is distributed under the Creative Commons Attribution 4.0 International licence, as recorded in the arXiv licence metadata for this exact version and shown on the abstract page https://arxiv.org/abs/2601.12846v1. A full rights notice is delivered at licenses/arxiv-2601.12846-CC-BY-4.0.txt. Characters: every retained excerpt is pure ASCII, so the delivered engine, XeLaTeX with its default Latin Modern text font, reports no missing character.
\begin{definition} \label{def:binding-circle}
Let $D$ be a link diagram.
A simple closed curve $\gamma$ is called a \textbf{binding circle for a three-page presentation} of $D$ if it satisfies the following conditions:
\begin{enumerate}
    \item The curve $\gamma$ intersects $D$ in finitely many points.
    \item All crossings of $D$ lie inside $\gamma$.
    \item Each arc of $D$ cut by $\gamma$ is of one of the following three types:
    \begin{enumerate}
        \item it lies outside $\gamma$;
        \item it lies inside $\gamma$ and participates only in over-crossings;
        \item it lies inside $\gamma$ and participates only in under-crossings.
    \end{enumerate}
    \item No two arcs of the same type are adjacent along $\gamma$.
\end{enumerate}
\end{definition}

\begin{lemma} \label{lem:binding-circle-existence}
Let $D$ be a link diagram with $n$ crossings.
Then there exists a binding circle $\gamma$ that induces a three-page presentation with at most $3n+1$ arcs.
\end{lemma}

\begin{lemma} \label{lem:extended}
Let $D$ be a non-split reduced link diagram with $n$ crossings, and let $Y$ be an extended spanning tree of $D$.
Then there exists a binding circle $\gamma$ that induces a three-page presentation of $D$ with at most $3n+1-m$ arcs,
where $m$ denotes the number of $2$-cells contained in $Y$.
\end{lemma}

\begin{proof}
Let $D$ be a non-split reduced link diagram with $n$ crossings, and let $Y \subset X(D)$ be an extended spanning tree.
Write $m$ for the number of $2$-cells contained in $Y$.
We construct a binding circle $\gamma$ as the boundary of a sufficiently small regular neighborhood of $Y$ in $S^2$.

Since $Y$ is connected and contractible, a regular neighborhood $\mathcal{N}(Y)$ is a topological disk.
Hence its boundary $\partial \mathcal{N}(Y)$ is a simple closed curve.
After a sufficiently small perturbation, we may assume that $\gamma:=\partial \mathcal{N}(Y)$ is transverse to $X^{(1)}(D)$ and meets each $1$-cell in finitely many points.

We now estimate the number of intersection points between $\gamma$ and the diagram $D$.
Each $0$-cell (crossing) of $X(D)$ has valence $4$ in the $1$-skeleton, so there are $4n$ local half-edges in total.
If $\gamma$ is taken as in Lemma~\ref{lem:binding-circle-existence} from a spanning tree, then one obtains $4n-(n-1)=3n+1$ intersections.
We explain how the presence of $2$-cells in $Y$ decreases this count.

Choose a spanning tree $T$ of the $1$-skeleton $Y^{(1)}$.
Because $Y$ is connected and contractible and contains all $0$-cells, we may view $T$ as a spanning tree of the underlying graph of $D$.
Let $\gamma_0$ be the binding circle obtained as the boundary of a sufficiently small regular neighborhood of $T$.
As in Lemma~\ref{lem:binding-circle-existence}, after a perturbation we may assume that $\gamma_0$ intersects every edge of $T$ exactly once and intersects each edge not in $T$ at most once.
In particular, $\gamma_0$ yields a three-page presentation with at most $3n+1$ arcs.

We now modify $\gamma_0$ using the $2$-cells contained in $Y$.
Let $f$ be a $2$-cell of $Y$.
By the definition of an extended spanning tree, no two $2$-cells of $Y$ share a common $1$-cell.
Hence we may perform the following local modification independently for different $2$-cells.

Consider a small neighborhood of $f$.
The boundary $\partial f$ is a cycle in the $1$-skeleton $X^{(1)}(D)$.
Since $f \subset Y$, the regular neighborhood $\mathcal{N}(Y)$ contains a neighborhood of $f$,
so the curve $\gamma=\partial \mathcal{N}(Y)$ runs outside $\partial f$ and, after a small perturbation, intersects exactly once the unique $1$-cell of $\partial f$ that was previously crossed by $\gamma_0$ in order to pass between the two sides of $\partial f$.
Pushing $\gamma_0$ across $f$ replaces this intersection by a segment of $\gamma$ lying in the interior of $f$ and therefore removes exactly one intersection point with the diagram.
Because distinct $2$-cells in $Y$ do not share a $1$-cell, these $m$ local modifications remove $m$ distinct intersection points.

After performing this modification for every $2$-cell in $Y$, we obtain a simple closed curve $\gamma=\partial \mathcal{N}(Y)$ such that
$$
|\gamma \cap D| \le (3n+1) - m.
$$
By construction, $\gamma$ is the boundary of a regular neighborhood of a connected, contractible subcomplex, so it satisfies the first three conditions in Definition~\ref{def:binding-circle}.
Moreover, since $Y$ contains no pair of $2$-cells sharing a common $1$-cell, a sufficiently small perturbation ensures that $\gamma$ meets every $1$-cell of $X(D)$, and the adjacency condition in Definition~\ref{def:binding-circle} follows exactly as in the proof of Lemma~\ref{lem:binding-circle-existence} (alternating case), with the same local avoidance modification in the non-alternating case.

Therefore, $\gamma$ is a binding circle inducing a three-page presentation of $D$ with at most $3n+1-m$ arcs.
\end{proof}

\end{document}
